\section{Related Work}
\label{sec:related}

Prior work provides the ingredients for deployment reliability
decisions but not their combination: reliability engineering
contributes the vocabulary, agent benchmarks the trace setting, and
probabilistic verification the model-based reasoning. The missing piece
is an audited fit from LLM-agent traces---a model with uncertainty,
diagnostics, and shared semantics for common metrics.

%% [REDUCTION] The full-width method-comparison table is suppressed for the
%% 12-page body; the three paragraphs below plus the positioning summary
%% state the same taxonomy. Retained in source for the artifact.
\iffalse
\begin{table*}[!t]
\centering
\caption{Comparison of \textsc{TraceToChain} with related evaluation
and verification frameworks.}
\label{tab:method-comparison}
\scriptsize
\setlength{\tabcolsep}{3pt}
\renewcommand{\arraystretch}{1.12}
\begin{tabular}{@{}>{\raggedright\arraybackslash}p{0.22\textwidth}
                >{\raggedright\arraybackslash}p{0.15\textwidth}
                >{\raggedright\arraybackslash}p{0.2\textwidth}
                >{\raggedright\arraybackslash}p{0.22\textwidth}
                >{\raggedright\arraybackslash}p{0.1\textwidth}@{}}
\toprule
\textbf{Approach} & \textbf{Input} & \textbf{Reliability semantics} &
\textbf{Diagnostics / uncertainty} & \textbf{LLM-agent fit} \\
\midrule
\mbox{Classical SRE / NHPP~\cite{cheung1980,goelokt}} &
Aggregate failures &
Growth / failure intensity &
Aggregate intervals, no trace states &
Indirect \\
\tstriped
\mbox{Agent benchmark metrics~\cite{chen2021codex,rdc2025}} &
Outcomes / traces &
Success, pass$@k$, pass$^k$, RDC &
Metric-specific, limited fit checks &
Direct / split \\
\mbox{PRISM / Storm~\cite{prism2011,storm2017}} &
User model &
PCTL queries &
Verification, no trace fitting &
Indirect \\
\tstriped
\mbox{TriCEGAR / ProbGuard~\cite{tricegar2026,pro2guard2025}} &
Empirical traces &
Learned Markov abstraction &
Refinement diagnostics &
Related \\
\thighlight
\mbox{\textbf{\textsc{TraceToChain} (this work)}} &
\textbf{LLM-agent traces} &
\textbf{First-passage DTMC reliability} &
\textbf{AIC/KS diagnostics, posterior/bootstrap intervals} &
\textbf{Direct} \\
\bottomrule
\end{tabular}
\end{table*}
\fi


\paragraph{Classical SRE.}
Classical software reliability gives this paper its vocabulary and
analytic tools. The dependability taxonomy of Avi{\v z}ienis et al.\
separates reliability, availability, safety, and failure
semantics~\cite{avizienis2004dependability}, while SRE
practice~\cite{beyer2016sre} turns such concepts into measurable
service objectives. State-based modeling is equally central:
Cheung~\cite{cheung1980} used Markov chains for program-level software
reliability, Trivedi and Bobbio~\cite{trivedi2017reliability} give a
broader engineering treatment, and Kemeny and Snell~\cite{kemenysnell}
developed the absorbing-chain framework and fundamental matrix
$N=(I-Q)^{-1}$ used by our closed forms. Reliability-growth models add
a complementary aggregate view~\cite{goelokt,musa1987}. None of these
is by itself a trace-level evaluation method for LLM agents;
Proposition~\ref{thm:nhpp} connects them by recovering the NHPP form as
a rare-failure scaling limit of the fitted agent chain.

\paragraph{Agent benchmarks.}
Agent benchmarks show why trace-level reliability semantics are
needed. ReAct~\cite{yao2023react}, Toolformer~\cite{schick2023toolformer},
and Reflexion~\cite{shinn2023reflexion} established agent
trajectories that interleave reasoning, tool use, and revision steps;
ToolBench~\cite{qin2023toolllm} and $\tau$-bench~\cite{taub2024} make
API choice, tool errors, and tool-agent-user interaction part of the
evaluation surface. SWE-bench~\cite{swebench2024},
WebArena~\cite{zhou2023webarena}, ALFWorld~\cite{shridhar2021alfworld},
and Voyager~\cite{wang2023voyager} move evaluation further toward
software, web, embodied, and open-ended environments where success is
long-horizon and context-dependent. MAST~\cite{mast2025} diagnoses
multi-agent failure modes, and AgentBench~\cite{agentbench2024}
broadens the interactive evaluation space beyond one task family.
Together, these benchmarks expose the trace setting; they do not by
themselves unify the reported metrics. pass$@k$~\cite{chen2021codex}
originated as a repeated-sampling success metric for code generation,
and agent evaluation later adapted it alongside the reliability decay
curve (RDC)~\cite{rdc2025}, which tracks how repeated-trial reliability
\emph{degrades as task duration grows}. Our $d$-step reliability $R(d)$
(Definition~\ref{def:reliab}) is the complementary step-budget view: a
non-decreasing cumulative success probability indexed by the step
budget rather than a duration-indexed decay. We retain the RDC name for
continuity and state the modelled quantity explicitly in
\S\ref{sec:unify}.

\paragraph{Formal methods for agents.}
Formal methods supply the model-based reasoning layer. Probabilistic
computation tree logic (PCTL)~\cite{hansson1994pctl} supports
reasoning about probabilistic time and reliability properties, and
model-checking texts~\cite{baier2008principles} place such logics in a
broader verification framework. Given a probabilistic model,
PRISM~\cite{prism2011} and Storm~\cite{storm2017} can verify PCTL
properties. More recent trace-driven frameworks, including
TriCEGAR~\cite{tricegar2026} and ProbGuard~\cite{pro2guard2025}, learn
DTMC or Markov decision process (MDP) abstractions from empirical
traces. We therefore do not claim trace-to-DTMC abstraction as the
main theoretical novelty. Our focus is its reliability-engineering use
for LLM agents: \textsc{TraceToChain} fits an absorbing DTMC with a
composite AIC$\,\wedge\,$KS goodness-of-fit certificate and uncertainty
quantification, and SS9 (\S\ref{sec:ss9}) evaluates held-out recovery on
controlled MAST-style corpora using KS distance and
$L_\infty^{\mathrm{RDC}}$ error. The first-passage view explains why
pass$@k$, pass$^k$, and the RDC are restrictions of a single
distribution (Theorem~\ref{thm:unify}), while the Goel--Okumoto
NHPP appears as a rare-failure scaling limit of the fitted chain
(Proposition~\ref{thm:nhpp}).

\noindent\textbf{Positioning summary.}
These traditions are complementary rather than competing: classical
reliability models provide semantics, agent benchmarks provide traces,
and probabilistic-verification tools provide model-based queries, but
none of them fits and audits a reliability model from LLM-agent traces.
\textsc{TraceToChain}'s distinction is the
audited combination---traces define the estimate, diagnostics test the
abstraction, uncertainty bounds the reported quantities, and
first-passage semantics connect horizon reliability, perturbations, and
common LLM-agent metrics---which yields a reliability argument from
trace data rather than another scalar score.

%% ================================================================
